\documentclass[11pt]{article}
\usepackage[margin=1in]{geometry}
\usepackage{amsmath, amssymb, amsthm}
\usepackage{booktabs}
\usepackage{enumitem}
\usepackage[hidelinks]{hyperref}

\newcommand{\code}[1]{\texttt{#1}}
\newcommand{\SVD}{\mathrm{SVD}}
\DeclareMathOperator*{\argmin}{arg\,min}

\title{Choosing the factor count when the pool is completed by EM}
\author{Cross-validation on held-out observed entries}
\date{September 14, 2026}

\begin{document}
\maketitle

\section{Why the usual criteria cannot be used here}

Every factor-count decision so far rests on a PCA of the complete-case pool:
twelve columns, rank at most eleven, each row's $r$ loadings fitted with
$12-r$ degrees of freedom. That is the regime in which extra factors behaved
like noise directions and we clamped the count to five or six.

Once the loadings are estimated by EM on all $N \approx 400$ dataset-periods,
that constraint is gone --- but so is the ability to read the count off a
spectrum. The completion returns
\begin{equation}
  \hat X(r) \;=\; \argmin_{\mathrm{rank}(Z)\le r}
  \bigl\| \mathcal{P}_\Omega(P - Z) \bigr\|_F^2 ,
\end{equation}
whose missing entries are filled at rank $r$ by construction. Its singular
values therefore satisfy $s_{r+1}(\hat X(r)) \approx 0$: the spectrum shows a
cliff exactly where we put one. Running Bai--Ng, Ahn--Horenstein or a scree rule
on $\hat X(r)$ recovers the assumption, not the data. The same circularity
defeats any criterion applied after completion.

What \emph{is} available, and is not circular, is the completion's ability to
predict entries it never saw.

\section{The criterion}

\subsection{Held-out entries}

Partition the observed set $\Omega$ at random into a training and a validation
part,
\begin{equation}
  \Omega \;=\; \Omega_{\mathrm{tr}} \,\dot\cup\, \Omega_{\mathrm{va}},
  \qquad
  \Pr\bigl[(i,t) \in \Omega_{\mathrm{va}}\bigr] = \pi \quad \text{i.i.d.},
  \qquad \pi = 0.1 .
\end{equation}
Crucially the split is over \emph{observed} cells only: entries missing in the
data stay missing in both. Run the completion on the training cells alone,
\begin{equation}
  \hat Z(r) \;=\; \argmin_{\mathrm{rank}(Z)\le r}
  \bigl\| \mathcal{P}_{\Omega_{\mathrm{tr}}}(P - Z) \bigr\|_F^2 ,
\end{equation}
by the same impute-and-refit iteration used in estimation (treating
$\Omega_{\mathrm{va}}$ as missing), and score it on the held-out cells:
\begin{equation}
  \mathrm{CV}(r) \;=\;
  \Biggl( \frac{1}{|\Omega_{\mathrm{va}}|}
  \sum_{(i,t)\in\Omega_{\mathrm{va}}} \bigl( P_{it} - \hat Z(r)_{it} \bigr)^2
  \Biggr)^{1/2} .
  \label{eq:cv}
\end{equation}
Average over $B$ independent splits to remove the split noise, and report
\begin{equation}
  \overline{\mathrm{CV}}(r) = \frac{1}{B}\sum_{b=1}^{B}\mathrm{CV}_b(r),
  \qquad
  \mathrm{se}(r) = \frac{1}{\sqrt{B}}\,\mathrm{sd}_b\bigl(\mathrm{CV}_b(r)\bigr).
\end{equation}

\subsection{What the curve means}

$\overline{\mathrm{CV}}(r)$ falls while additional components capture structure
that generalizes across cells, and rises once they begin fitting cell-specific
noise. The minimizer,
\begin{equation}
  \hat r = \argmin_r \overline{\mathrm{CV}}(r),
\end{equation}
is the rank the pool supports \emph{given its missingness pattern} --- which is
the quantity we have never been able to estimate, because the complete-case
spectrum answers a different question (how many directions twelve columns of one
dataset exhibit).

Two conventions are worth reporting alongside the minimizer. The
one-standard-error rule
\begin{equation}
  \hat r_{1\mathrm{se}} = \min\bigl\{ r : \overline{\mathrm{CV}}(r) \le
  \overline{\mathrm{CV}}(\hat r) + \mathrm{se}(\hat r) \bigr\}
\end{equation}
picks the most parsimonious count statistically indistinguishable from the best,
which matches how we have been treating factor counts elsewhere. And the
relative improvement
\begin{equation}
  \delta(r) = \frac{\overline{\mathrm{CV}}(r-1) - \overline{\mathrm{CV}}(r)}
                   {\overline{\mathrm{CV}}(r-1)}
\end{equation}
shows where the curve flattens, which is usually more informative than the exact
argmin when the minimum is shallow.

\section{Block missingness: what the criterion can and cannot see}

The validation cells are drawn from $\Omega$, so they are cells that
\emph{some} dataset-period observes. Consequently:

\begin{itemize}[leftmargin=*]
  \item The criterion measures how well rank-$r$ structure predicts
        \emph{within-block} variation --- a copula coefficient of the
        PSID-replica at a date whose other copula coefficients are seen, a
        wealth coefficient of the SCF-replica, and so on. This is precisely the
        information EM adds relative to complete cases, so it is the right thing
        to validate.
  \item It cannot validate the model's extrapolation into cells that are
        \emph{never} observed --- consumption-and-wealth cross terms outside the
        one dataset that measures both. No held-out sample exists there, because
        no data exist there. $\hat r$ should therefore be read as ``the rank the
        observed cells support'', with the understanding that the completion's
        behaviour in never-observed blocks rests on the factor structure rather than
        on evidence.
\end{itemize}

A refinement that partially addresses the second point: hold out entire
\emph{columns} rather than scattered cells (a group-$k$-fold over
dataset-periods). Then $\mathrm{CV}(r)$ measures how well the loadings estimated
without a given period predict that period's observed coefficients, which is
closer to the use the model makes of them. Scattered-cell CV is the standard
choice and the cheaper one; column CV is a useful robustness check.

\section{Interaction with what the count is used for}

$\hat r$ answers a measurement question: how many directions the coefficient
pool supports. It does not answer the dynamic question the state space cares
about --- whether those directions are \emph{persistent} enough, or aligned with
the conditional statistics the paper reports. Two consequences.

First, $\hat r$ is an upper bound worth respecting rather than a target to hit:
if the curve is flat from $r=6$ to $r=10$, the one-standard-error count is the
right choice, because the extra factors cost
$O\bigl((r+q)^2\bigr)$ parameters in the innovation correlation block for a
reconstruction gain the criterion says is negligible.

Second, the previous evidence against large $r$ was collected in the
complete-case regime and does not transfer: factors 6--10 were then estimated
from twelve columns, which is why they behaved like noise. Under EM they are
estimated from $N \approx 400$, and under
\code{persistence\_prior = :per\_factor} an unneeded factor can settle at low
persistence with small innovations instead of injecting jitter. The empirical
question of whether that is what actually happens is settled by the run, not by
the criterion.

\section{Cost and implementation}

Each evaluation of $\mathrm{CV}(r)$ is one EM completion: $O(\mathrm{iters})$
truncated SVDs of an $n \times N$ matrix with $n \approx 1730$, $N \approx 400$,
$r \le 12$. With a randomized or truncated SVD this is seconds; over a grid of
ten ranks and $B = 5$ splits it is a few minutes at prep time, with no MCMC. The
proposed interface is a diagnostic switch that prints
\begin{center}
\begin{tabular}{rrrr}
\toprule
$r$ & $\overline{\mathrm{CV}}(r)$ & $\mathrm{se}$ & $\delta(r)$ \\
\midrule
3 & \dots & \dots & --- \\
4 & \dots & \dots & \dots \\
\multicolumn{4}{c}{\dots} \\
\bottomrule
\end{tabular}
\end{center}
together with $\hat r$ and $\hat r_{1\mathrm{se}}$, and exits before the sampler
so the choice is made by inspection rather than automatically.

\paragraph{Reference.} The estimator being validated is the rank-constrained
completion of Section 2; the cross-validation scheme is the standard one for
matrix completion (hold out observed entries, refit, score the held-out set),
adapted only in that the missingness here is structural rather than random, which
is what Section 3 qualifies.

\end{document}
